% !TEX program = xelatex
\documentclass{英语试卷}
\begin{document}
\examtitle{唐山市 2023—2024 学年度高三年级第一学期期末考试}{英语}
\par 唐山市 2023—2024 学年度高三年级第一学期期末考试\par
\par 英语试题卷\par
\begin{notice}
\par 1. 答卷前，考生务必将自己的姓名、准考证号填写在答题卡上。\par
\par 2. 回答选择题时，选出每小题答案后，用铅笔把答题卡上对应题目的答案标号涂黑。如需改动，用橡皮擦干净后，再选涂其他答案标号。回答非选择题时，将答案写在答题卡上，写在本试卷上无效。\par
\par 3. 考试结束后，将本试卷和答题卡一并交回。\par
\end{notice}
\parttitle{第一部分 听力（共两节，满分30分）}
\par 做题时，先将答案标在试卷上，录音内容结束后，你将有两分钟的时间将试卷上的答案转涂到答题卡上。\par
\sectitle{第一节（共5小题：每小题1.5分，满分7.5分）}
\par 听下面 5 段对话。每段对话后有一个小题，从题中所给的 A、B、C 三个选项中选出最佳选项。听完每段对话后，你都有 10 秒钟的时间来回答有关小题和阅读下一小题。每段对话仅读一遍。\par
\question{1. What will the woman do tonight?}
\begin{choices}(3)
  \item Enjoy a concert.
  \item Have a new dish.
  \item Make a plan for her sister.
\end{choices}
\question{2. What will the speakers probably order?}
\begin{choices}(3)
  \item Soup.
  \item A meat pie.
  \item Chicken.
\end{choices}
\question{3. How did the man know about Focus?}
\begin{choices}(3)
  \item On the Internet.
  \item From its member.
  \item Through a relative.
\end{choices}
\question{4. What makes the man a little sad?}
\begin{choices}(3)
  \item A small gathering.
  \item Tom's departure.
  \item A new neighbor's arrival.
\end{choices}
\question{5. What are the two speakers talking about?}
\begin{choices}(3)
  \item A local team.
  \item A weekend plan.
  \item A football match.
\end{choices}
\sectitle{第二节（共15小题：每小题1.5分，满分22.5分）}
\par 听下面 5 段对话或独白。每段对话或独白后有几个小题，从题中所给的 A、B、C 三个选项中选出最佳选项。听每段对话或独白前，你将有时间阅读各个小题，每小题 5 秒钟；听完后，各小题将给出 5 秒钟的作答时间。每段对话或独白读两遍。\par
\question{6. How many days did the trip take?}
\begin{choices}(3)
  \item 21.
  \item 24.
  \item 60.
\end{choices}
\question{7. What does the woman think of the landscape in Mongolia?}
\begin{choices}(3)
  \item It's boring.
  \item It's exciting.
  \item It's frightening.
\end{choices}
\par 听第 7 段材料，回答第 8 至 9 题。 微信公众号 IAI English\par
\question{8. Why does the man complain?}
\begin{choices}(3)
  \item The staff said bad words.
  \item The cooling system failed.
  \item The seat was uncomfortable.
\end{choices}
\question{9. How would the man describe the solution?}
\begin{choices}(3)
  \item Disappointing.
  \item Acceptable.
  \item Perfect.
\end{choices}
\par 听第 8 段材料，回答第 10 至 12 题，\par
\question{10. Where will the woman probably go?}
\begin{choices}(3)
  \item Peru.
  \item India.
  \item Spain.
\end{choices}
\question{11. What food will the woman avoid?}
\begin{choices}(3)
  \item Meat.
  \item Peanuts.
  \item Cheese.
\end{choices}
\question{12. When will the interview take place?}
\begin{choices}(3)
  \item On Friday.
  \item On Thursday.
  \item On Tuesday.
\end{choices}
\question{13. What will the woman do in the interview?}
\begin{choices}(3)
  \item Present slides.
  \item Display equipment.
  \item Share her experience.
\end{choices}
\par 听第 9 段材料，回答第 13 至 16 题。\par
\question{14. What did the man dream of. doing?}
\begin{choices}(3)
  \item Going to Kenya.
  \item Becoming an athlete.
  \item Studying in university.
\end{choices}
\question{15. Why did the man regret?}
\begin{choices}(3)
  \item He did not practice earlier.
  \item He did not train hard enough.
  \item He did not follow his parents' advice.
\end{choices}
\question{16. How does the man feel after studying the course?}
\begin{choices}(3)
  \item Relieved.
  \item Regretful.
  \item Curious.
\end{choices}
\question{17. What's the possible relationship between the two speakers?}
\begin{choices}(3)
  \item Mother and son.
  \item Fellow classmates.
  \item Teacher and student.
\end{choices}
\par 听第 10 段材料，回答第 17 至 20 题。\par
\question{18. Why did Chandi start the expedition?}
\begin{choices}(3)
  \item To make history.
  \item To do some research.
  \item To seek more challenges.
\end{choices}
\question{19. What helped Chandi survive the depressing days?}
\begin{choices}(3)
  \item Self-reflection.
  \item Inspiring voices.
  \item Adventurous stories.
\end{choices}
\question{20. Which word can best describe Chandi?}
\begin{choices}(3)
  \item Determined.
  \item Competitive.
  \item Considerate.
\end{choices}
\parttitle{第二部分 阅读理解（共两节，满分 50 分）}
\sectitle{第一节（共15小题：每小题2.5分，满分37.5分）}
\par 阅读下列短文，从每题所给的 A、B、C 和 D 四个选项中，选出最佳答案。\par
\passagelabel{A}
\par PUBLIC SPEAKING\par
\par INSTRUCTOR: Dr. Elise Morrison\par
\par OBJECTIVES:\par
\par Public Speaking is designed to improve your understanding of the communication process and to develop your skills as a public speaker and as a critic of public speech.\par
\par ASSIGNMENTS:\par
\par Throughout the course you will deliver four prepared speeches of increasing length and complexity. You will prepare outlines for each speech. You will also write two short papers that analyze various aspects of oral presentation and rhetoric.\par
\par Speaking: Speech of Introduction: Introduce yourself in 2 minutes: The goal is for your audience to remember your name and learn something about you.\par
\par Speech to Teach: The 4-minute speech will teach us about an idea or how to do something. The emphasis is\par
\par on organization as well as on establishing common ground and credibility.\par
\par Speech to Initiate an Action: You have 6 minutes to persuade the audience to take a specific action. The focus is on motivation and follow-through.\par
\par Speech to Change an-Attitude: You have 8-10 minutes to change your audience's attitudes on a topic that you will have pre-tested for agreement/ disagreement. The emphasis is on argumentation, credibility, and persuasion.\par
\par Interviews: You will receive in-class coaching on how to persuasively communicate your skills in a job interview.\par
\par Written: Rhetorical Analysis of a Famous Speech: Write an essay that analyzes the rhetorical strategies used in a famous speech of your choosing.\par
\par Self-Assessment/ Course Evaluation: The purpose of this assignment is to assess your work in the course and the course itself.\par
\par GRADING\par
\par Oral Presentations and Written Components 90\%\par
\par Rhetorical Analysis of Famous Speech (20\%)\par
\par Speech to Initiate an Action(20\%) Outline(10\%) Speech(10\%)\par
\par Speech to Change an Attitude (30\%) Outline(15\%) Speech(15\%) Self-evaluation(5\%)\par
\par Class Participation 10\%\par
\question{21. Where is this text probably taken from?}
\begin{choices}(2)
  \item A speech announcement.
  \item A lesson design.
  \item An academic article.
  \item A course guide.
\end{choices}
\question{22. Which speech focuses on the structure?}
\begin{choices}(2)
  \item Speech to Teach.
  \item Speech of Introduction.
  \item Speech to Initiate an Action.
  \item Speech to Change an Attitude.
\end{choices}
\question{23. How much does written work account for?}
\begin{choices}(2)
  \item 15\%.
  \item 25\%.
  \item 30\%.
  \item 35\%.
\end{choices}
\passagelabel{B}
\par Growing up in Louisiana, across the Mississippi River from New Orleans, Barry Gui lot loved exploring the wetlands with his friends when he was a child. However, as he grew older, a growing concern about the disappearance of wetlands took root. "Imagine if the New Orleans Saints, our football team, were playing on a field that was actually wetlands," Barry says. "By halftime, that football field would be gone."\par
\par In 1988, Barry became a middle school science teacher. With the intention of bringing home to his kids the importance of wetlands, he founded the LaBranche Wetland Watchers Service-Learning Project, "adopting" a small part of LaBranche Wetlands near their school.\par
\par Netland Watchers activities are tied to academic subjects. Water-quality monitoring, for example, teaches students to use graphs to compare data from different time periods-part of the math curriculum. After Hurricane Katrina, the salinity (amount of salt in the water) was four times as high as before. Students made graphs to show that. As part of English composition, they wrote about seeing jellyfish (水母), which had never come that far before because jellyfish live only in salty water.\par
\par "We obtained more than we would just sitting behind a desk with a book, because you're out there and you're getting all wet and muddy and having fun as you learn," says Kurt, a seventh grader.\par
\par Barry is very proud of his students. "It's amazing what middle school kids can accomplish when they get the chance," he says. One of his seventh graders wrote, "If the animals and plants could talk, I think they would say we're their heroes. That's the way I feel when we do our work in the wetlands."\par
\question{24. Why is "football field" mentioned in the first paragraph?}
\begin{choices}(2)
  \item To explain wetlands are limited.
  \item To express Barry's love of football.
  \item To stress the urgency of saving wetlands.
  \item To show Barry's worry about his football team
\end{choices}
\question{25. What is Paragraph 3 mainly about?}
\begin{choices}(2)
  \item The necessity of making graphs.
  \item The importance of hands-on learning.
  \item The connection of the activities to schooling.
  \item Relationship between teachers and students.
\end{choices}
\question{26. How does Kurt feel about the activities?}
\begin{choices}(2)
  \item They are messy.
  \item They are challenging.
  \item They are dull.
  \item They are rewarding.
\end{choices}
\question{27. What can be inferred about Barry's students?}
\begin{choices}(2)
  \item They will teach like Barry.
  \item They will ask for more chances.
  \item They're full of a sense of achievement.
  \item They're acknowledged as great heroes.
\end{choices}
\passagelabel{C}
\par Growing up in suburbia, yard work was punishment. I'd take almost anything else. I didn't like being outside, and I especially didn't like bugs or dirt.\par
\par When I left home at 17 with a backpack of belongings, I figured my gardening days were behind me. There were no yards to work in-only concrete! The concrete jungle that awaited me seemed like a welcome escape from the outdoor chores that had once felt like a sentence.\par
\par It wasn't until my mid-20s that I, thinking it was the adult thing to do, brought home a small daisy(雏菊) in a pot to put on my fire escape. Less than 24 hours later, a squirrel(松鼠) arrived and dragged it away, pot and all. That's New York for you! I took it as a sign and didn't try any more plants in that apartment.\par
\par When I married, my wife and I bought a house in outer Brooklyn. There was enough room to have a garden if we wanted one. We decided to see if we could grow anything, mostly for the novelty factor. Much to my surprise, digging in the dirt felt like fun and totally not a chore! In the early mornings, I would bend down to inspect the soil waiting for the plants to appear. Watering at dusk, I whispered to the baby plants, expecting them to keep growing. The plants kept the neighborhood squirrels fed and provided a place for bumblebees to nap. I felt a rush of power and excitement.\par
\par As the flowers and vegetables bloomed, I noticed something in me was growing as well. Since running away from home as a teenager, I've always done my best to build a shelter for myself wherever I was living. I hung pictures on the walls. I built new holiday traditions with friends. I painted the walls bright colors. But it always felt like maybe that deep connection to a particular space was missing. In gardening, I finally found this new layer of depth. I guess you could say I was putting down roots.\par
\par There is an inspiring, self-care practice in the time and energy it takes to plant, water, and wait and wait to see if the magic happens. It's the act of growing that I find the sweetest. What are you growing?\par
\question{28. How does the author interpret his departure at 17?}
\begin{choices}(2)
  \item As a necessary step for growth.
  \item As a great relief.
  \item As an escape from social expectations.
  \item As a bittersweet transition.
\end{choices}
\question{29. What can we learn about the author at his mid-20s?}
\begin{choices}(2)
  \item There was a shift in his perspective.
  \item He regarded a squirrel as a bad sign.
  \item He began to miss the days before 17.
  \item A sense of responsibility grew in him.
\end{choices}
\question{30. What did the author harvest from gardening?}
\begin{choices}(2)
  \item A colorful life.
  \item Friends' respect.
  \item Creative ideas.
  \item A sense of belonging.
\end{choices}
\question{31. The question at the end is intended to stimulate the readers to consider \blanknum{}.}
\begin{choices}(2)
  \item their own sources of growth
  \item the importance of gardening
  \item the output of the growing plants
  \item their favorite flowers and vegetables
\end{choices}
\passagelabel{D}
\par The Greek philosopher Heraclitus taught that you can't step into the same river twice, for you aren't the same person at each visit, and the water is ever flowing. It is a powerful way to represent the reality of impermanence: Everything is always changing.\par
\par Many people have tense relationships with change. They deny it, resist it or attempt to control it-the result of which is almost always some combination of stress, anxiety, and burnout. It doesn't have to be that way.\par
\par A concept called allostasis can help. Developed in the late 1980s by a neuroscientist, Peter Sterling, and a biologist, Joseph Eyer, allostasis is defined as "stability through change," elegantly showing the concept's double meaning: The way to stay stable through the process of change is by changing, at least to some extent.\par
\par The river of change has been flowing mercilessly, and it shows no signs of letting up. Societally, we've gone through a pandemic(大流行病), its economic decline and the widespread adoption of social media, the combination of which has shifted how we live and work. In our personal lives, we continue to do what we have always done: relocate, start jobs, quit jobs, change jobs, get promoted, experience illness, get married, have children, retire and on and on.\par
\par Finding ways in this river requires both ruggedness and flexibility. To be rugged is to be tough and determined, to know your core values and what you stand for. To be flexible is to actively respond to changing circumstances, to adapt and bend easily without breaking, to grow and even change your mind. Put these qualities together and the result is a great endurance, one that helps you maintain your strong core even in fragile moments. It allows you to step into allostasis's cycle of order, disorder and reorder - which is, of course, one and the same with stepping into Heraclits's river — and to chart it skillfully and whenever possible, to your own benefit.\par
\question{32. What do we know about allostasis?}
\begin{choices}(2)
  \item It is a contradictory concept.
  \item It tends to focus on change.
  \item It advocates dynamic stability.
  \item It claims absolute adaptability.
\end{choices}
\question{33. Which of the following best explains "letting up" underlined in paragraph 4?}
\begin{choices}(2)
  \item Being pumped.
  \item Slowing down.
  \item Getting polluted.
  \item Rising up.
\end{choices}
\question{34. What will you gain from practising allostasis according to the last paragraph?}
\begin{choices}(2)
  \item Popularity.
  \item Confidence.
  \item Recognition.
  \item Perseverance.
\end{choices}
\question{35. What is the best title for the passage?}
\begin{choices}(2)
  \item Say No to Stability
  \item Stop Resisting Change
  \item Fight for Flexibility
  \item Balance Life and Work
\end{choices}
\sectitle{第二节（共5小题；每小题2.5分，满分12.5分）}
\par 阅读下面短文，从短文后的选项中选出能填入空白处的最佳选项。选项中有两项为多余选项，\par
\par As the trees drop their leaves and frosts advance in, we prepare our homes and ourselves for the winter months. \blanknum{36} When the seasons change, we know that we have to change our behavior to keep ourselves safe and warm. Our feelings change like the seasons. Laughing. joyful and positive like summer; sad, angry, and hopeless like winter. \blanknum{37} However, it would not be practical to feel wonderful all the time.\par
\par The phrase "emotional wintering" was popularized by Katharine May. She mentions emotional wintering asks us to see our difficult or uncomfortable feelings as winters which we can get ready for and live through. \blanknum{38} Sometimes the emotional winters are caused by an event that makes us feel as though we want to withdraw into ourselves. \blanknum{39} We need to take time to replenish (补充) them while accepting that we are sad. Whatever the cause, emotional wintering encourages us to accept that our emotional seasons are normal and we should manage them as best we can.\par
\par \blanknum{40} To manage through challenging times, we need to adjust our habits, reassess our expectations of ourselves, and ensure that we sustain ourselves appropriately. This process gives us the space to explore the full range of emotions that make us human, and allows us to safely sit with a feeling, in relative comfort, until it passes.\par
\par With the management, we believe we will be well by the time the next summer rolls around.\par
\begin{choices}(1)
  \item It prepares us for the coming spring.
  \item Different emotions come to all of us.
  \item Emotional wintering asks us to change.
  \item We can grow from challenging experience.
  \item Thick socks and cosy blankets are brought out.
  \item We expect to always live in our emotional summers.
  \item At other times we simply use up our supplies of positivity.
\end{choices}
\parttitle{第三部分 语言知识运用（共两节，满分 30 分）}
\sectitle{第一节（共15小题；每小题1分，满分15分）}
\par 阅读下面短文，从短文后各题所给的 A、B、C 和 D 四个选项中，选出可以填入空白处的最佳选项。\par
\par Dolphins. That's how I got into trouble. My family were eating dinner when I \blanknum{41} it up. I mentioned a friend posting overly inspirational quotes accompanied by photos of rainbows, sunsets, heart-shaped clouds, you \blanknum{42} it. "Now she's moved on to dolphins, I wonder if killer whales are \blanknum{43}." Everyone in my family \blanknum{44}, except my 17-year-old daughter, Maggie. "You're snarky(爱挖苦的), Mom. You make fun of everything." I wanted to argue, but Maggie gave me the teenage \blanknum{45} shoulder.\par
\par All night I thought about what Maggie had said. Was she \blanknum{46}? Was I too snarky? Actually, I wasn't a \blanknum{47} person. But around people I was comfortable with, I did tend to \blanknum{48} things.\par
\par So for this New Year's Resolution, I \blanknum{49} I'd stop being snarky. The first few days, I wasn't expressing the snark, but the \blanknum{50} kept running in my head. I, however, controlled them. Gradually, I managed not only to stop being snarky, but to be \blanknum{51} to others.\par
\par Several days later, I was in the regular supermarket. The check-out clerk was a sour woman who rarely said anything apart from the \blanknum{52} "welcome". Noticing she seemed sad, I didn't give the \blanknum{53} in my usual way. Instead, I just \blanknum{54} for her. I did something I'd never done before. I smiled at her. She smiled back. I left the store feeling \blanknum{55} somehow.\par
\cz{41}{tore}{put}{brought}{picked}
\cz{42}{like}{name}{draw}{share}
\cz{43}{alive}{next}{dangerous}{past}
\cz{44}{screamed}{rose}{froze}{laughed}
\cz{45}{cold}{broad}{bony}{tense}
\cz{46}{pessimistic}{reliable}{right}{guilty}
\cz{47}{greedy}{ridiculous}{rigid}{mean}
\cz{48}{set aside}{put up with}{make fun of}{take down}
\cz{49}{decided}{agreed}{pretended}{guessed}
\cz{50}{arguments}{comments}{jokes}{plans}
\cz{51}{kind}{blind}{humble}{plain}
\cz{52}{controlled}{required}{interrupted}{teased}
\cz{53}{assistance}{information}{response}{suggestion}
\cz{54}{felt}{behaved}{arose}{fought}
\cz{55}{clearer}{tougher}{milder}{lighter}
\sectitle{第二节（共10小题；每小题1.5分，满分15分）}
\par 阅读下面短文，在空白处填入 1 个适当的单词或括号内单词的正确形式，\par
\par Chinese classical literature holds a vital position in the cultural fabric of China. \blanknum{56} (root) in a history extending for thousands of years, these literary treasures convey the core of Chinese thought, providing a window into the cultural transition of one of the world's oldest \blanknum{57} (civilization).\par
\par Chinese classical literature is like a treasure chest filled with \blanknum{58} (wise), good advice, and deep thoughts. The Confucian classics, such as The Analects and Mencius, are like the important building blocks that have helped create the rules for \blanknum{59} people should behave in China for a very long time. The poems, particularly those by Tang poets like Li Bai, are like timeless songs, through \blanknum{60} the poets express the Chinese feelings towards life, nature, and their intense emotions.\par
\par Classical literature is like a large history book that narrates the transition of dynasties. Books like Records of the Grand Historian and Romance of the Three Kingdoms are not just stories, but special windows that let us see what many people \blanknum{61} experience) during those times.\par
\par In a \blanknum{62} (rapid) changing world, Chinese classical literature is like a cultural cornerstone, \blanknum{63} (encourage)a sense of connection and identity for the Chinese people. The special thing about it is that it doesn't get old. This makes \blanknum{64} a great medium for different generations to feel a strong link to the values and dreams that make \blanknum{65} the Chinese way of thinking.\par
\parttitle{第四部分 写作（共两节，满分 40 分）}
\sectitle{第一节（满分15分）}
\par 假定你是李华。你校即将举办“科技与学习”的主题活动，计划邀请姊妹学校 Highgate School 进行线上互动。请你写一封邮件，沟通相关事宜。\par
\par 内容包括:\par
\par （1）活动的时间、目的；\par
\par (2) 连线的时间、话题。\par
\par 注意:\par
\par (1) 写作词数应为 80 左右;\par
\par (2) 请按如下格式在答题纸的相应位置作答。\par
\par Dear Sir or Madam,\par
\par Li Hua\par
\sectitle{第二节（满分 25 分）}
\par 阅读下面材料，根据其内容和所给段落开头语续写两段，使之构成一篇完整的短文。\par
\par My darling grandma has a collection of photo albums kept over several decades. I enjoyed joking with grandma about her preferred old-fashioned method of photo storage, but I know that these albums are precious for her. They are the records of our family and our identity.\par
\par One day while examining one of grandma's albums, I came across photos of my mother and her brother as children in Disney World in the 1980's. Grandma looked thoughtful and smiled. "Your mother and your uncle had such an amazing time," she said. "I'm glad I was able to make that vacation happen for them. Every child should visit Disney at least once and just get to be a kid... happy and carefree." As I continued examining the photos, I noticed something which didn't sit right with me. In every photo my grandmother was in, she was smiling in an unsettled, half-hearted way where her smile didn't quite reach her eyes. When I asked her about it, she gently closed the album and said, "Oh, now enough of these old photos. Let's talk about what we're going to do for Christmas next month."\par
\par That evening, I quietly asked my mother about the photos and the family trip to Disneyland. Mom breathed out. She explained that my grandparents were constantly fighting and that they had been having problems in their marriage. "I think that was probably the last trip we took together as a family before they separated," Mom said. "Your grandma was so sad the entire week we were there, but she did everything she could to make sure your uncle and I had a great experience. We had been begging to go to Disney for so long. She organized the entire trip for our pleasure. I'm sure it was all physically and emotionally tiring for her."\par
\par That night I turned restlessly in my bed thinking about the devotion my grandmother had for her children and about how even today she continues to put the needs of everyone in our family above her own. Grandma does so much for her children and grandchildren every day. 注意：\par
\par (1) 续写词数应为 150 左右;\par
\par (2) 请按如下格式在答题卡的相应位置作答。\par
\writingpara{Paragraph 1: Shouldn't we do something for her?}
\writinglines{8}
\writingpara{Paragraph 2: We pulled over at Disney World, which surprised my grandma.}
\writinglines{8}
\end{document}
